\section{Правила имитационной модели}
\label{sec:model}

\subsection{Инициализация}
\code{init_simulation()} проверяет параметры и сбрасывает генератор по
\code{seed}, статистику, очередь и историю предыдущего эксперимента.
Затем создаётся каталог, загружаются или генерируются партии и учитывается
закупочная стоимость всего начального запаса. До создания постоянных
покупателей выполняется уценка на день 0. Поэтому уценённые лекарства уже
влияют на выбор их регулярных покупок.

Параметры каталога для лекарства с индексом $i$:
\begin{itemize}
  \item дозировка $25(1+\lfloor i/6\rfloor)$ мг;
  \item форма с индексом $i\bmod 6$, группа с индексом $i\bmod 5$;
  \item оптовая цена --- равномерное целое от 50 до 600 рублей;
  \item срок годности --- равномерное целое от 45 до 180 дней;
  \item минимальный запас --- равномерное целое от 10 до 25 упаковок.
\end{itemize}
Формы: таблетки, суспензия, спрей, мазь, капсулы и раствор. Группы:
сердечно-сосудистые, антибиотики, эндокринные, обезболивающие и
противоаллергические. Эти названия --- данные учебной модели.

Без CSV каждое лекарство получает $C$ упаковок, где $C$ --- параметр
\code{initial_count}:
первая партия содержит $\lfloor C/2\rfloor$ и случайный срок от 1 до срока
годности лекарства; вторая --- оставшиеся упаковки и полный срок годности.
При $C=0$ партии не создаются. \code{initial_packages} считается до первого дня.

\subsection{Порядок одного дня}
Вызовы выполняются строго последовательно:
\begin{enumerate}
  \item \code{manage_warehouse(day)} устанавливает день, списывает истёкшие
        партии, применяет уценку, собирает потери и принимает готовые поставки.
  \item \code{generate_daily_events(day)} добавляет плановые и разовые заказы
        в общую очередь.
  \item \code{process_daily_order(day)} доставляет подготовленное вчера,
        считает нагрузку, выделяет товар из очереди на завтра и создаёт
        заявки пополнения.
  \item \code{cleanup_random_customers()} сохраняет нужных покупателей,
        обновляет индексы очереди и публикует результат дня.
\end{enumerate}

Движок контролирует фазу и номер дня. Начинать необходимо с дня 1;
пропуски, повторное начало дня, обработка заказов до генерации и повторное
завершение дня приводят к \code{std::logic_error}.
Период заканчивается после дня $N$: дополнительные дни для завершения
доставок не выполняются. В первый день доставок нет, а подготовленное в день
$N$ остаётся в \code{pending_deliveries}.

Результат из \code{simulation_result()} обновляется при инициализации и
завершении дня. Чтение между фазами не даёт отдельного согласованного снимка
промежуточных изменений.

\subsection{Поток заказов}
Обычный спрос моделируется 180 независимыми попытками в день с вероятностью
\begin{equation}
  p(m)=0{,}75\exp\left(-\frac{m}{60}\right), \qquad
  X_d\sim\operatorname{Binomial}(180,p(m)), \qquad
  \mathbb{E}[X_d]=180p(m).
\end{equation}
Число курьеров напрямую в эту вероятность не входит. При повышении наценки
ожидаемое число обычных заказов уменьшается. Плановые заказы постоянных
покупателей добавляются к этому потоку.

Каждый обычный заказ содержит от 1 до 4 различных лекарств и от 1 до 4
упаковок каждой позиции. Вес выбора лекарства равен 3, если на складе есть
его уценённая непустая партия, действительная на текущий день; иначе вес равен 1.
После выбора позиция исключается из дальнейшего выбора в том же заказе.
Лекарства без доступного запаса также могут попасть в запрос с обычным весом.
Вес относится к лекарству целиком, а не к отдельной партии.

Для постоянного покупателя период равномерно выбирается от 3 до 7 дней,
а первый день покупки --- от 1 до величины периода. Состав регулярной покупки
генерируется один раз. Если \code{next_purchase_day <= day}, создаётся один
плановый заказ и следующий день сдвигается на период до значения больше
текущего дня. Плановые покупки формируются перед обычным потоком.

У покупателей случайно определяется наличие карты. ФИО, телефоны и адреса
синтетические. Подготовленный и выполненный заказ хранит собственную копию
покупателя; удаление разовых клиентов из глобального списка не меняет историю.

\subsection{Выделение товара и очередь}
Очередь обрабатывается в порядке поступления. Сначала доставляются вчерашние
покупки, затем выделяется товар для не более чем $15M$ новых покупок.
Склад уменьшается именно при подготовке. Для каждого запроса партии
рассматриваются по возрастанию срока годности; при одинаковом сроке
сохраняется исходный порядок.

При текущем дне $d$ для доставки завтра подходит партия только при $e\geq d+1$.
Из подходящих партий выдаётся доступное количество, не превышающее запрос.
Частично выполненный заказ считается одной подготовленной покупкой;
невыданный остаток его запроса повторно в очередь не помещается.
Если не выдано ни одной упаковки, заказ удаляется из очереди и увеличивает
\code{empty_orders}; отдельной доставки для него нет.

Лимит $15M$ относится к успешным подготовленным покупкам: пустые заказы
не занимают места. Когда лимит достигнут, ещё не рассмотренные заказы
остаются в очереди без выделенного товара. После освобождения вместимости
они будут обработаны в последующие дни. Поэтому доставка выполняется
на следующий день после \emph{подготовки}, а исходный заказ может ждать дольше.

\code{Order::items} сохраняет исходный запрос. Фактическая выдача доступна
отдельно в \code{supplied}; цена рассчитывается только по ней.
Поле \code{prepared_day} фиксирует день снятия товара со склада,
\code{ordered_day} --- день поступления, \code{delivered_day} --- день доставки.

\subsection{Курьеры}
Для покупки с порядковым номером $i$ в списке вчерашней подготовки
назначается индекс
\begin{equation}
  c_i=(i+d-1)\bmod M, \qquad i=0,1,\ldots,D_d-1.
\end{equation}
Распределение отличается между курьерами не более чем на одну покупку.
Движок повторно считает нагрузку по реальным выполненным заказам.
Максимум составляет 15 покупок на курьера.

Если $7M\leq D_d\leq15M$, каждый курьер получает от 7 до 15 покупок.
При $D_d<7M$ выполнить нижнюю границу для всех невозможно;
\code{courier_limits_satisfied} получает значение \code{false}.
Фиктивные заказы не создаются. \code{late_deliveries} считает доставки,
выполненные позже дня, следующего за \code{ordered_day}.
В текстовом выводе курьеры нумеруются с 1; C++ и JSON используют индексы с 0.

\subsection{Уценка и списание}
Партия годна в свой \code{expiration_day}. Списание происходит при $e<d$;
в потери поступает закупочная стоимость оставшихся упаковок, а партия удаляется.
При $0\leq e-d\leq30$ продажная база уменьшается вдвое один раз.
Повторный вызов уценки в следующие дни не делит цену повторно.

Закупочная цена сохраняется в \code{BatchAccounting}. Поэтому
\code{StoreBatch::batch_price} после уценки описывает базу продажи, а не
себестоимость. Уценка уменьшает будущую выручку; сама по себе она не создаёт
запись в \code{write_off_losses}. Пустые партии удаляются при обработке склада.

Например, партия с $e=40$ уценивается на день 10, годна на день 40 и списывается
на день 41. Для доставки на день 40 её можно выделить на день 39;
на день 40 для завтрашней доставки она уже не подходит.

\subsection{Пополнение склада}
После подготовки заказов \code{SupplyManager} суммирует остатки действительных
партий каждого лекарства. Если $q_i<s_i$ и ожидающей заявки для лекарства нет,
создаётся заявка объёмом
\begin{equation}
  Q_i=3s_i-q_i.
\end{equation}
Условие строгое: при $q_i=s_i$ заявка не создаётся. Ожидающие поставки
не добавляются к фактическому остатку, но блокируют повторную заявку на
то же лекарство.

Задержка равномерно выбирается из 1, 2 и 3 дней. Поставка добавляет новую
партию с оптовой ценой из текущего каталога. Её последний день годности
равен текущему дню $d$ плюс \code{shelf_life} лекарства.
Это правило действует и после задания начальных цен через CSV.
Закупочные расходы начисляются при фактическом поступлении;
стоимость ещё ожидающих поставок в расходы не входит.

\subsection{Цена и скидки}
Пусть $q_b$ --- количество выделенных упаковок из партии $b$, $v_b$ ---
её текущая продажная база. Сумма после наценки, до скидок:
\begin{equation}
  S=\sum_b q_bv_b\left(1+\frac{m}{100}\right).
\end{equation}
Основная скидка определяется так:
\begin{equation}
  \delta_0=
  \begin{cases}
    \delta_{\mathrm{card}}, & \text{есть карта},\\
    \delta_{\mathrm{bulk}}, & \text{карты нет и }S>1000,\\
    0, & \text{иначе}.
  \end{cases}
\end{equation}
Постоянному клиенту добавляется $\delta_{\mathrm{regular}}$, после чего
применяется общий предел:
\begin{equation}
  \delta=\min\bigl(\delta_0+mathbf{1}_{\mathrm{regular}}
                   \delta_{\mathrm{regular}},\delta_{\max}\bigr),
  \qquad P=S\left(1-\frac{\delta}{100}\right),\quad 0\leq\delta_{\max}\leq9.
\end{equation}
Скидка по карте и скидка крупной покупки одновременно не складываются.
При $S=1000$ скидки крупной покупки нет. Порог сравнивается с суммой уже
после уценки и наценки, но до скидок покупателю.

При базовой цене 100 рублей, 10 упаковках и наценке 25\% сумма $S=1250$.
Постоянный клиент с картой и настройками по умолчанию получает
$\min(5+5,9)=9$\%, то есть платит 1137{,}50 рубля. Если партия уценена,
$S=625$ и цена того же клиента составляет 568{,}75 рубля; закупочная
стоимость этих упаковок остаётся 1000 рублей.

\subsection{Статистика, прибыль и балансы}
В дневной \code{totals} и в JSON-полях дня хранятся накопленные суммы
с начала эксперимента. Выручка признаётся при доставке, а не при подготовке.
Начальный запас считается приобретённым перед началом эксперимента.

Обозначим: $I$ --- выручка доставленных покупок, $B$ --- все закупки,
$A$ --- закупочная стоимость доступного запаса, $R$ --- стоимость выделенного
товара до доставки, $C_D$ --- себестоимость доставленного, $L$ --- потери.
Используются соотношения:
\begin{align}
  B &= A+R+C_D+L,\\
  \Pi &= I-B+A+R = I-C_D-L,\\
  F &= I-B.
\end{align}
Здесь $\Pi$ --- прибыль модели, $F$ --- денежный поток. Потери нельзя повторно
вычитать из выражения $I-B+A+R$: они уже отражены в стоимости исчезнувшего
запаса. Ожидающая доставка учитывается в $R$ и пока не создаёт выручку.
Стоимость запаса $A$ и $R$ определяется по закупочным ценам, даже после уценки.

Движок сравнивает две формулы прибыли при инициализации и после каждого дня.
Допуск составляет $10^{-7}\max(1,B)$; превышение вызывает
\code{std::logic_error}. Существующие тесты также проверяют количественный
баланс упаковок и соответствие выручки списку реальных доставок.

Деньги хранятся как \code{double}; специального округления каждой операции
до копеек нет. Консоль печатает денежные итоги с двумя знаками, JSON использует
17 значащих цифр. Возможные длинные дробные представления в JSON отражают
арифметику типа \code{double}.
